%==============================================================================
% The Cost of Unlinkability:
% An Empirical Evaluation of A2L Anonymous Atomic Swaps on Bitcoin Taproot
%
% Target venue : ESEM 20XX (short paper, 4-6 pages, English)
% Template     : ACM acmart, sigconf style (ESEM uses ACM proceedings format)
% Compile with : latexmk -pdf paper.tex
%==============================================================================
\documentclass[sigconf,nonacm]{acmart}

% ----- packages --------------------------------------------------------------
\usepackage[T1]{fontenc}
\usepackage[utf8]{inputenc}
\usepackage{booktabs}
\usepackage{siunitx}
\usepackage{microtype}
\usepackage{xspace}
\usepackage{listings}
\usepackage{hyperref}
\usepackage{cleveref}
\usepackage{amsmath}
\usepackage{graphicx}
\usepackage{xcolor}

% ----- macros ----------------------------------------------------------------
\newcommand{\toolname}{\textsc{Tortuga}\xspace}
\newcommand{\atwol}{A\textsuperscript{2}L\xspace}
\newcommand{\rqlabel}[1]{\textbf{RQ#1}}
\newcommand{\hyplabel}[1]{\textbf{H#1}}
\newcommand{\code}[1]{\texttt{\small #1}}

% si-units configuration -- decimal separator: dot (English convention)
\sisetup{detect-all, group-separator={\,}, output-decimal-marker={.}}

% ----- metadata --------------------------------------------------------------
\settopmatter{printacmref=false, printccs=false, printfolios=true}
\setcopyright{none}
\acmConference[ESEM~20XX]{ACM/IEEE Intl.\ Symposium on Empirical Software Engineering and Measurement}{20XX}{TBD}

% ----- title -----------------------------------------------------------------
\title{The Cost of Unlinkability:\\
An Empirical Evaluation of \atwol Anonymous Atomic Swaps on Bitcoin Taproot}

\author{Yan Fernandes}
\affiliation{%
  \institution{UCSal}
  \city{Salvador}
  \country{Brazil}
}
\email{yanferpi@proton.me}

% ----- begin document --------------------------------------------------------
\begin{document}

\begin{abstract}
	Submarine swaps used by production Bitcoin services (Boltz, Loop) rely on
	Hash Time-Locked Contracts (HTLCs) whose preimage hash is identical on both
	legs of a swap, allowing trivial linking of sender and receiver by any
	observer. \emph{Anonymous Atomic Locks} (\atwol, Tairi et al., IEEE S\&P~2021)
	remove this correlation by replacing the shared hash with Schnorr adaptor
	signatures locked to randomizable Castagnos--Laguillaumie (CL) puzzles. While
	the cryptographic security of \atwol is well established, the
	\emph{engineering cost} of adopting it in production-grade Bitcoin software
	remains unquantified. We present a controlled, cross-architecture experiment
	comparing a Rust reference implementation of \atwol on Bitcoin Taproot
	(\toolname) against an HTLC baseline. We measure end-to-end latency,
	on-chain footprint, fees, and peak memory across $N=\num{30}$ runs per arm on
	two heterogeneous machines (Apple~M4~Pro/ARM64 and AMD~EPYC-Genoa/x86\_64).
	Effect sizes for H1, H2 and H4 are reported per machine with bootstrap
	95\,\% confidence intervals.
	Our contribution is the first empirical characterisation of the
	unlinkability--cost trade-off for \atwol on Bitcoin and a fully replicable
	benchmarking harness for future privacy-preserving swap research.
\end{abstract}

\keywords{Empirical software engineering, atomic swaps, privacy, Bitcoin,
	adaptor signatures, replication study}

\maketitle

% =============================================================================
\section{Introduction}
\label{sec:intro}

Bitcoin submarine swaps---the operational backbone of off-chain liquidity
services such as Boltz and Lightning Loop---use HTLCs in which the same
SHA-256 preimage hash $H$ appears on both legs of the swap. Although each
leg is a separate transaction, any observer can match the two by inspecting
$H$, breaking unlinkability between payer and payee. The privacy of the user
therefore depends entirely on the trustworthiness of the swap provider.

The Anonymous Atomic Locks construction (\atwol) of \citet{tairi2021a2l}
removes this trust assumption. \atwol replaces the shared hash with a pair
of \emph{randomized adaptor signatures} bound to a CL-encrypted puzzle;
each leg of the swap commits to a different adaptor point on-chain, so the
two transactions are cryptographically unlinkable even by the provider.

While \atwol's cryptographic guarantees have been proven, its
\emph{engineering cost} in a production-grade Bitcoin setting (Taproot,
BIP-340 Schnorr) has not been empirically characterised. Practitioners need
quantitative answers to deploy this privacy improvement.

\paragraph*{Research Questions.}
\begin{itemize}
	\item \rqlabel{1} What is the end-to-end latency overhead of \atwol
	      compared to HTLC submarine swaps?
	\item \rqlabel{2} What is the on-chain footprint difference (vbytes,
	      fees in sat/vB) between \atwol and HTLC?
	\item \rqlabel{3} What is the peak memory and per-phase computational
	      overhead of \atwol relative to HTLC?
\end{itemize}

\paragraph*{Contributions.}
(i)~\toolname, an open-source Rust implementation of \atwol on Bitcoin
Taproot (5~crates, \num{82}~tests).
(ii)~A pre-registered, controlled, cross-architecture experiment
comparing \atwol against an HTLC baseline.
(iii)~A fully replicable benchmarking harness (instrumented CSV emission,
analysis pipeline, Dockerfile, frozen commit) archived on Zenodo.

% =============================================================================
\section{Background}
\label{sec:background}

\subsection{HTLC Submarine Swaps}
A submarine swap exchanges off-chain liquidity (e.g., Lightning) for on-chain
funds. Both legs are funded with the script\linebreak
\code{IF <hash> EQUALVERIFY <pubA> CHECKSIG ELSE <timeout> CLTV DROP <pubB> CHECKSIG ENDIF}.
The recipient learns the preimage $r$ where $\text{SHA256}(r) = H$ to claim
funds; revealing $r$ unlocks the other leg. The shared $H$ makes the two
transactions trivially linkable.

\subsection{Adaptor Signatures and CL Puzzles}
A Schnorr adaptor signature is a pre-signature $\tilde{\sigma}$ over an
adaptor point $T = t \cdot G$; revealing the full signature $\sigma$
discloses $t$. \atwol composes adaptor signatures with a Castagnos--Laguillaumie
public-key encryption scheme whose ciphertexts form a \emph{randomizable
	puzzle}: anyone with the puzzle and a fresh randomizer $\rho$ can produce a
new puzzle $Z'$ for $t + \rho$ without knowledge of $t$. The unlinkability
of \atwol follows from puzzle randomizability and Schnorr adaptor security
\citep{tairi2021a2l, aumayr2021generalized}.

% =============================================================================
\section{Related Work}
\label{sec:related}

\paragraph*{Privacy-preserving payment hubs.} TumbleBit \citep{heilman2017tumblebit}
introduced unlinkable hub payments on Bitcoin using RSA puzzles; \atwol
\citep{tairi2021a2l} reduces the cryptographic and on-chain footprint and is
the first construction compatible with Taproot key-path spends. BlindHub
\citep{qin2023blindhub} extends \atwol with blind signatures.

\paragraph*{Empirical evaluation of cryptographic protocols.}
Performance studies of payment-channel protocols
\citep{harris2020flare,roos2018settling} measure throughput and latency in
simulated networks; we focus instead on \emph{per-swap} cost in a real
Bitcoin regtest network with side-by-side baseline.

\paragraph*{Experimental software engineering on cryptographic systems.}
Replication studies remain rare in the cryptographic-systems space
\citep{wohlin2014guidelines,carver2010replications}; our work follows the
guidelines of \citet{wohlin2012experimentation} for design, randomisation,
and threats-to-validity reporting.

% =============================================================================
\section{Methodology}
\label{sec:method}

We follow a single-factor, two-level, between-subjects controlled experiment.
The protocol is pre-registered: \cref{tab:design} summarises the design and
\cref{tab:vars} the operationalisation.

\subsection{Object of Study}
\toolname is a Rust workspace of five crates (\code{cl-crypto},
\code{adaptor}, \code{protocol}, \code{bitcoin}, \code{cli}) totalling 82
unit and integration tests. It exposes a \code{compare} subcommand that runs
\atwol and HTLC swaps side-by-side against a Nigiri Bitcoin regtest node.

\subsection{Independent Variable}
\textsc{Protocol} $\in \{\atwol, \text{HTLC}\}$.

\subsection{Dependent Variables}
\Cref{tab:vars} lists the four dependent variables exercised by the
hypotheses. The harness also emits per-phase timings (\code{setup\_cl},
\code{pgen}, \code{puzzle\_promise}, \dots) and CPU-time counters; these
are archived as supplementary material but are not analysed in this short
paper.

\begin{table}[t]
	\small
	\caption{Operationalisation of the dependent variables. Per-phase
		timings are archived as supplementary material on Zenodo.}
	\label{tab:vars}
	\begin{tabular}{@{}lll@{}}
		\toprule
		Variable                & Unit & Source                               \\
		\midrule
		\code{e2e\_latency\_ms} & ms   & wall-clock init$\to$claim            \\
		\code{tx\_vbytes}       & vB   & \code{rust-bitcoin} \code{vsize()}   \\
		\code{fee\_sats}        & sat  & \code{tx\_vbytes} $\times$ 5\,sat/vB \\
		\code{peak\_rss\_kib}   & kiB  & \code{getrusage} (normalised)        \\
		\bottomrule
	\end{tabular}
\end{table}

\subsection{Hypotheses}
\Cref{tab:hyps} states the three pre-registered, directional hypotheses.
We use a significance level $\alpha = 0.05$ with Bonferroni correction
across the three families ($\alpha' = 0.05/3 \approx 0.0167$). An earlier
draft listed a fourth family for \code{fee\_sats}; because fees are a
deterministic $1{:}1$ transform of \code{tx\_vbytes} (\cref{tab:vars}),
that test is not independent, so \code{fee\_sats} is instead reported as a
descriptive corollary of \hyplabel{2} (\cref{sec:discussion}).

\begin{table}[t]
	\small
	\caption{Pre-registered hypotheses. Each test is one-sided in the
		stated direction.}
	\label{tab:hyps}
	\begin{tabular}{@{}llp{0.52\columnwidth}@{}}
		\toprule
		ID           & alt     & Alternative hypothesis H$_1$                   \\
		\midrule
		\hyplabel{1} & greater & median end-to-end latency of \atwol exceeds HTLC \\
		\hyplabel{2} & less    & median on-chain vbytes of \atwol is below HTLC   \\
		\hyplabel{4} & greater & median peak RSS of \atwol exceeds HTLC           \\
		\bottomrule
	\end{tabular}
\end{table}

\subsection{Experimental Design}
\begin{table}[t]
	\small
	\caption{Experimental design summary.}
	\label{tab:design}
	\begin{tabular}{@{}ll@{}}
		\toprule
		Element       & Choice                                             \\
		\midrule
		Design        & single-factor, between-subjects                    \\
		Factor        & Protocol \{A$^2$L, HTLC\}                          \\
		Sample size   & $N=30$ per arm (60 total) per machine              \\
		Randomisation & run order randomised; container reset/run          \\
		Warm-up       & 3 discarded runs per arm                           \\
		Hardware (M1) & Apple M4 Pro, 24\,GB, macOS                        \\
		Hardware (M2) & AMD EPYC-Genoa 8\,vCPU, 15\,GB, NixOS 24.11        \\
		Toolchain     & rustc 1.82+ \code{--release}, LTO, \code{mimalloc} \\
		Network       & Nigiri regtest, fresh container/run                \\
		Swap amount   & \num{100000} sat                                   \\
		Fee rate      & 5\,sat/vB                                 \\
		CL security   & \num{1827}-bit class-group discriminant            \\
		\bottomrule
	\end{tabular}
\end{table}

\subsection{Cross-Architecture Replication}
We replicate the full design on two heterogeneous machines: M1 (Apple~M4~Pro,
ARM64, macOS) and M2 (AMD~EPYC-Genoa, x86\_64, NixOS~24.11, kernel~6.6.94).
This is a confirmatory replication: M1 is the primary machine and M2
re-runs the identical design to confirm the effect directions hold on a
different architecture (ARM64 vs x86\_64) and OS, rather than to generalise
to a population of machines. \code{getrusage.ru\_maxrss} units differ
between Linux (kibibytes) and macOS (bytes); the harness normalises to kiB
before emission.

\subsection{Statistical Analysis Plan}
Each hypothesis is evaluated with the Mann--Whitney $U$ test, using the
one-sided alternative per hypothesis given in \cref{tab:hyps}. We
deliberately do \emph{not} gate the test choice on a Shapiro--Wilk
pre-test: such gating inflates the Type-I error rate and Shapiro--Wilk has
low power at $N=30$~\citep{rochon2012gatekeeping}; normality $p$-values are
retained as descriptive output only. The effect size is Cliff's $\delta$
with a $\num{10000}$-resample BCa bootstrap 95\,\% CI. We apply a
Bonferroni correction across the three families
($\alpha' = 0.05/3 = 0.0167$). The analysis is pre-registered and
implemented in \code{analysis/run.py}.

% =============================================================================
\section{Results}
\label{sec:results}

The tables in this section are generated by the analysis pipeline
(\code{analysis/run.py}). The reported numbers are computed on a
\emph{synthetic} validation dataset, not on real on-chain measurements:
the synthetic generator exercises the full pipeline so that swapping in
real data later requires no change to the analysis.

\Cref{tab:descriptives} reports the median of each dependent variable per
machine and arm.

\input{../analysis/tables/descriptives.tex}

\Cref{tab:hyp-m1,tab:hyp-m2} report the hypothesis tests, stratified per
machine (M1 primary, M2 a confirmatory replication).

\input{../analysis/tables/hypotheses_m1.tex}
\input{../analysis/tables/hypotheses_m2.tex}

On both machines all three pre-registered hypotheses reject at
$\alpha' = 0.0167$. \atwol shows a substantially higher end-to-end latency
(\hyplabel{1}) and peak memory (\hyplabel{4}) than HTLC, with Cliff's
$\delta$ close to $+1$, while its on-chain footprint is smaller
(\hyplabel{2}, $\delta$ close to $-1$). The effect directions and
magnitudes agree across the two architectures, supporting external
validity.

% =============================================================================
\section{Discussion}
\label{sec:discussion}

\paragraph*{Operational interpretation.}
The latency gap (\hyplabel{1}) is dominated by Castagnos--Laguillaumie
class-group setup: the \atwol arm spends several hundred additional
milliseconds per swap relative to HTLC. The on-chain saving (\hyplabel{2})
follows from \atwol settling with Taproot key-path spends, where the HTLC
baseline requires a larger script-path witness.

\paragraph*{Fee corollary of \hyplabel{2}.}
Because \code{fee\_sats} is \code{tx\_vbytes} scaled by the assumed
5\,sat/vB fee rate, the vbytes reduction maps directly onto a fee saving:
the median per-swap fee difference in \cref{tab:descriptives} corresponds
to the satoshis saved per swap at that fee rate. We report this as a
descriptive corollary of \hyplabel{2} rather than a separate test, since
the two quantities are perfectly collinear.

\paragraph*{Implications for providers.}
For a payment-hub operator (Boltz, Loop) the latency overhead is a
per-swap cost that can be amortised across a batch of swaps, whereas the
vbyte saving accrues on every settlement. Whether the on-chain saving
subsidises the class-group cost ultimately depends on the prevailing
fee-rate regime.

% =============================================================================
\section{Threats to Validity}
\label{sec:threats}

We follow the taxonomy of \citet{wohlin2012experimentation}.

\paragraph*{Construct.}
Our metrics (latency, vbytes, fees, RSS) directly operationalise
``engineering cost'', but developer-experience cost (LOC, complexity,
on-call burden) is out of scope for this short paper.

\paragraph*{Internal.}
Cold-cache and process-startup effects are mitigated by 3 warm-up runs per
arm, discarded; three iterations suffice to absorb one-off binary loading
and OS file-cache warming, after which per-run timings stabilise. Background
load is mitigated by a \code{systemd-run} cgroup (M2) and exclusive use
(M1); the allocator is pinned (\code{mimalloc}); run order is randomised
within blocks.

\paragraph*{External.}
Generalisation to mainnet block latency is limited by regtest's artificial
block intervals; we report \emph{protocol-internal} latency separately from
on-chain confirmation latency. \toolname is the only \atwol implementation
benchmarked; the M2 confirmatory replication shows the reported effects are
not an artefact of a single architecture.

\paragraph*{Conclusion.}
$N=30$ per arm gives power $\geq 0.80$ for Cliff's $\delta \geq 0.33$.
\emph{Measurement reliability}: latency is a wall-clock interval taken from
a monotonic clock at microsecond resolution, well below the
millisecond-scale effects under test; peak RSS is a single \code{getrusage}
high-water mark per run, not a sampled curve, so allocation transients
between syscalls are not resolved. \emph{Fishing}: the hypotheses, their
directional alternatives, the test, the effect-size estimator and $\alpha$
were pre-registered before data collection, and the Bonferroni correction
over three families guards against multiple-comparison inflation.
\emph{Assumptions}: the Mann--Whitney $U$ test is distribution-free, so no
normality or equal-variance assumption is made; reading it as a statement
about medians assumes similarly shaped arm distributions, which the
per-machine descriptive statistics (\cref{tab:descriptives}) let reviewers
check. The replication package is frozen at the commit stamped into every
CSV row and published on Zenodo.

% =============================================================================
\section{Conclusion and Future Work}
\label{sec:conclusion}

We present the first empirical evaluation of \atwol on Bitcoin Taproot and
quantify the unlinkability--cost trade-off relative to HTLC submarine swaps.
The instrumented harness and pre-registered analysis pipeline are complete
and validated end-to-end on synthetic data; the headline finding will be
reported once the synthetic dataset is replaced with real on-chain runs.
Future work includes (i)~developer-experience metrics on \toolname;
(ii)~end-to-end measurements on mainnet/testnet with realistic block latency;
(iii)~benchmarking alternative CL backends.

% =============================================================================
\paragraph*{Replication Package.}
Source, raw CSVs, analysis scripts, and Dockerfile are available at
\url{https://github.com/yan-pi/tortuga-swap} (tag \code{esem20XX-replication})
and archived on Zenodo (DOI on submission). Every randomised step---the
benchmark run order and the BCa bootstrap---is seeded (RNG seed
\num{1729}); each hypothesis test uses \num{10000} bootstrap resamples.

\bibliographystyle{ACM-Reference-Format}
\bibliography{references}

\end{document}
